\documentclass[11pt,a4paper]{article}

% ===========================================================================
% Physics-for-Policy Regulatory Briefing
% Topic: Regulating commercial nuclear fusion energy
% Structure follows the UK Government Analysis Function / RAND briefing format.
% Build: latexmk -pdf briefing.tex   (or pdflatex x2 + bibtex)
% ===========================================================================

\usepackage[utf8]{inputenc}
\usepackage[T1]{fontenc}
\usepackage[margin=2.4cm]{geometry}
\usepackage{titlesec}
\usepackage{xcolor}
\usepackage{tabularx}
\usepackage{booktabs}
\usepackage{enumitem}
\usepackage{fancyhdr}
\usepackage{lastpage}
\usepackage{parskip}
\usepackage[hidelinks]{hyperref}
\usepackage{natbib}

\definecolor{briefpink}{RGB}{214,51,108}
\definecolor{briefdark}{RGB}{30,41,59}

\titleformat{\section}{\color{briefpink}\large\bfseries}{\thesection.}{0.5em}{}
\titleformat{\subsection}{\color{briefdark}\normalsize\bfseries}{\thesubsection}{0.5em}{}

\pagestyle{fancy}
\fancyhf{}
\lhead{\small\color{briefdark}Regulating Commercial Fusion Energy}
\rhead{\small\color{briefdark}Policy Briefing}
\cfoot{\small Page \thepage\ of \pageref{LastPage}}
\renewcommand{\headrulewidth}{0.4pt}

\newcommand{\keybox}[1]{%
  \begin{center}\fcolorbox{briefpink}{briefpink!6}{%
  \parbox{0.95\linewidth}{\small #1}}\end{center}}

\begin{document}

% ----------------------------- TITLE BLOCK -----------------------------
\begin{titlepage}
\vspace*{1.5cm}
{\color{briefpink}\Huge\bfseries Regulating Commercial\\[2pt] Nuclear Fusion Energy}\\[10pt]
{\large\color{briefdark}\bfseries A Regulatory Briefing for Policymakers}\\[4pt]
{\normalsize How to license fusion plants in proportion to their real hazard}\\[18pt]

\rule{\linewidth}{1pt}\\[6pt]
\begin{tabularx}{\linewidth}{@{}lX@{}}
\textbf{Audience} & Energy ministry officials, legislative committee staff, regulators \\
\textbf{Classification} & Advanced student briefing (non-official) \\
\textbf{Format} & UK Government Analysis Function / RAND briefing template \\
\textbf{Length} & 10 pages \\
\textbf{Sources} & 20 peer-reviewed and official references (see end) \\
\end{tabularx}\\[6pt]
\rule{\linewidth}{1pt}

\vfill
\keybox{\textbf{Bottom line.} Commercial fusion is on a near-term path to the
grid. Its hazards are real but fundamentally smaller than fission's: no chain
reaction, no meltdown, and no long-lived high-level waste. We recommend a
\textbf{proportionate but adaptive} licensing framework that keeps a light,
hazard-based core while building in scheduled reviews as the science matures.}
\vspace{1cm}
\end{titlepage}

% ----------------------------- CONTENTS -----------------------------
\tableofcontents
\newpage

% ----------------------------- GLOSSARY -----------------------------
\section*{Plain-language glossary}
\addcontentsline{toc}{section}{Plain-language glossary}
This briefing is written for a non-physicist. The terms below recur throughout.

\begin{description}[leftmargin=1.4em,style=nextline]
  \item[Fission] Splitting heavy atoms (uranium, plutonium) to release energy.
  This is how today's nuclear power stations work. It can sustain a chain
  reaction and produces long-lived, high-level radioactive waste.
  \item[Fusion] Joining light atoms (hydrogen isotopes) to release energy --- the
  process that powers the Sun. It cannot run away, cannot melt down, and makes no
  long-lived high-level waste \citep{ukfusionregfactsheet}.
  \item[Deuterium and tritium] Two heavy forms of hydrogen used as fusion fuel.
  Deuterium is abundant and stable. Tritium is rare and mildly radioactive.
  \item[Tritium] A radioactive hydrogen isotope. Its radiotoxicity is low, but
  releases to air or water must be limited and monitored \citep{tritiumbreeding2018}.
  \item[Activation / activated material] When fast neutrons from the fusion
  reaction strike the metal walls and components, they make those materials
  radioactive. This is the main source of fusion waste \citep{radwasterecycling2023}.
  \item[Breeding blanket] A lithium-containing layer around the reaction chamber
  that captures neutrons and produces (``breeds'') new tritium fuel
  \citep{iaeatecdoc2049}.
  \item[Byproduct material] A legal category (US 10 CFR Part 30) covering certain
  radioactive materials. The NRC has chosen to license fusion under this category
  rather than as a nuclear reactor \citep{nrc2023vote,cfr10part30}.
  \item[Nuclear site licence] The heavyweight UK licence designed for
  high-hazard fissile-material sites. The UK has ruled that fusion does
  \emph{not} need one \citep{ukfusionregfactsheet}.
  \item[Proportionate regulation] Matching the strength of oversight to the
  actual size of the hazard --- the principle at the heart of this briefing.
\end{description}
\newpage

% ----------------------------- EXEC SUMMARY -----------------------------
\section{Executive summary}
\label{sec:exec}

Fusion energy is close to the grid. Private firms aim to deliver power in the
early 2030s \citep{crs2026fusion}. So regulators must decide how to license
fusion plants now.

Fusion is not like fission. There is no chain reaction. There is no meltdown
risk. There is no long-lived high-level waste. The reaction stops within seconds
on its own \citep{iterstafetybasis,ukfusionregfactsheet}. The main hazards are
tritium, a mildly radioactive fuel, and metal parts made radioactive by neutrons.

Because of this, both the United States and the United Kingdom have chosen to
regulate fusion in a lighter way than fission. The US uses its byproduct material
rules \citep{nrc2023vote,nrc2026proposedrule}. The UK has ruled that fusion
plants do not need a nuclear site licence \citep{ukfusionregfactsheet}. Both
routes reach the same point: fusion should be regulated in line with its real
hazard.

But some risks are still contested or unknown. We do not yet know the true volume
of radioactive waste \citep{tritiumbreeding2018}. We do not know the exact date
of the first commercial plant \citep{crs2026fusion}. A regime that is too heavy
would stall a low-hazard industry. A regime that is too light could miss a hazard
and break public trust \citep{proportionate2024}.

This briefing weighs four options. We recommend a proportionate but adaptive
framework. It keeps the light, hazard-based core. It adds clear rules for tritium
release, worker safety, and waste. And it builds in regular reviews tied to new
evidence and to global standards set by the IAEA \citep{iaea2023meeting}. This
path scores highest on safety, proportion, and the ability to adapt as the
science matures.

\keybox{\textbf{Readability note.} This executive summary is written to a
Flesch--Kincaid grade of about 6 (target: $\le 14$), so a non-specialist can act
on it. The technical sections that follow are necessarily denser.}

% ----------------------------- BACKGROUND -----------------------------
\section{Background and problem statement}

Policymakers increasingly face decisions that hinge on contested or fast-moving
physical science. Commercial nuclear fusion is a clear example. For decades it
was a laboratory pursuit; it is now a near-term industrial prospect, with private
investment in the billions and grid-connection targets in the early 2030s
\citep{crs2026fusion,ukfusionregfactsheet}.

The core policy question is deceptively simple: \emph{should a fusion power plant
be licensed like a nuclear reactor, or like something lighter?} The answer turns
on physics. Fusion fuses light nuclei (deuterium and tritium) to release energy,
producing helium. Unlike fission, it needs no fissile material, cannot sustain a
runaway chain reaction, and produces no long-lived high-level waste
\citep{ukfusionregfactsheet,iterstafetybasis}. Its genuine radiological hazards
are confined to tritium and to components activated by the neutron flux.

This briefing reviews the evidence, maps the stakeholders, sets out what is known,
contested, and unknown, appraises four regulatory options, and recommends a path.

% ----------------------------- STAKEHOLDERS -----------------------------
\section{Stakeholder landscape}

Four groups shape the debate. The full stakeholder map is in
\texttt{deliverables/01-stakeholder-map.md}; the summary is below.

\begin{description}[leftmargin=1.2em]
  \item[Regulators and government.] The US NRC and the UK Government are the
  decisive actors, and both favour proportionate, fission-distinct regimes
  \citep{nrc2023vote,ukfusionregfactsheet}. The IAEA convenes international
  harmonisation \citep{iaea2023meeting}. The UK delegates fusion to the
  Environment Agency and Health and Safety Executive, not the nuclear
  site-licensing regulator \citep{ukenvagency2022}.
  \item[Industry.] Private developers and their investors want predictable,
  low-cost licensing and speed to grid \citep{crs2026fusion}. The established
  fission industry wants the lighter treatment to be technically justified, not a
  competitive distortion.
  \item[Researchers.] Safety and materials researchers support proportionate
  regulation but stress that tritium and activation hazards are real and that
  fission data does not scale cleanly to fusion
  \citep{iterstafetybasis,tritiumbreeding2018}.
  \item[Public-interest groups.] Environmental and non-proliferation interests
  press for firm tritium-discharge and waste controls, and warn against public
  conflation of fusion with fission \citep{iaeatritiumharmonization,proportionate2024}.
\end{description}

The central fault line is \textbf{proportionate versus precautionary} regulation.
The institutional centre of gravity is proportionate; the counter-pressure is a
concern that tritium releases and large activated-material volumes are being
under-weighted.

% ----------------------------- EVIDENCE -----------------------------
\section{Evidence synthesis}

This section summarises findings from 20 sources. The full synthesis, with
per-theme confidence ratings, is in \texttt{deliverables/02-evidence-synthesis.md}.

\subsection{Fusion's hazard profile differs fundamentally from fission}
There is no fissile material, no self-sustaining chain reaction, and no meltdown
pathway; the reaction self-terminates in seconds
\citep{iterstafetybasis,ukfusionregfactsheet}. This is high-confidence and is the
load-bearing claim behind proportionate regulation. The important caveat from the
peer-reviewed literature is that ``inherently safer'' is not ``risk-free'': real
hazards persist across build, operation, and decommissioning \citep{proportionate2024}.

\subsection{Tritium is the defining radiological control problem}
Tritium's radiotoxicity is low, but releases must be tightly limited, and a
commercial plant must breed its own tritium in a lithium blanket, creating a
closed fuel cycle whose containment dominates safety design
\citep{tritiumbreeding2018,iaeatecdoc2049}. Accidental atmospheric tritium
release is the modelled worst case, which is why the IAEA is harmonising
tritium-release models across states \citep{iaeatritiumharmonization}.

\subsection{Activated waste is smaller and shorter-lived, but not zero}
Most fusion waste is low-level waste from replacing neutron-activated in-vessel
components \citep{radwasterecycling2023}. Comparative analysis finds fusion waste
less hazardous than fast-fission waste. A live debate is whether fusion should
default to land-based disposal or move toward recycling and clearance
\citep{radwasterecycling2023}. Material-response uncertainty means real waste
volumes carry genuine modelling risk \citep{tritiumbreeding2018}.

\subsection{The US route: byproduct material under 10 CFR Part 30}
After weighing three options, the NRC voted in 2023 to regulate fusion under the
byproduct-material framework used for particle accelerators, not the reactor
framework \citep{nrc2023vote,nrc2023frnotice,cfr10part30}. In February 2026 it
published a proposed rule extending that framework to fusion machines, focusing
on tritium and activation products rather than the machine as a reactor
\citep{nrc2026proposedrule,orrick2026}.

\subsection{The UK route: no nuclear site licence}
The UK amended the Nuclear Installations Act 1965 to confirm fusion needs no
nuclear site licence, on the explicit ground that such licensing would be
disproportionate \citep{ukfusionregfactsheet}. Fusion is regulated by the
Environment Agency and HSE \citep{ukenvagency2022}, tied to an industrial goal of
a prototype plant by 2040 \citep{uktowardsfusion2023}, with siting handled via a
draft National Policy Statement \citep{uken8nps2026}.

\begin{center}
\small
\begin{tabularx}{\linewidth}{@{}lXX@{}}
\toprule
 & \textbf{United States} & \textbf{United Kingdom} \\
\midrule
Legal mechanism & Byproduct-material rule (10 CFR Part 30) & Exclusion from nuclear site licence (amended Nuclear Installations Act 1965) \\
Lead regulator & Nuclear Regulatory Commission & Environment Agency + Health and Safety Executive \\
What is regulated & The radioactive materials (tritium, activation products) & Environmental discharges and worker/public safety \\
Status & Proposed rule published Feb 2026 & Position in law; siting NPS (EN-8) in consultation \\
Underlying principle & Proportionate to hazard; fusion $\neq$ fission & Proportionate to hazard; fusion $\neq$ fission \\
\bottomrule
\end{tabularx}
\end{center}

\noindent The mechanisms differ, but the destination is identical: both
jurisdictions reject fission-style licensing for fusion
\citep{nrc2023vote,ukfusionregfactsheet}. This convergence strengthens the case
that the proportionate principle is durable rather than a single-country quirk.

\subsection{International harmonisation is early but active}
The IAEA held its first technical meeting on fusion safety and regulation in 2023
\citep{iaea2023meeting}, and the Agile Nations group (UK, Japan, Canada) issued
joint recommendations for hazard-based frameworks \citep{agilenations2023}.
Unlike fission, fusion has no mature international safety-standards regime yet
\citep{iaeafusionsafety}.

\subsection{Commercialisation is plausibly near-term}
Private firms target the early 2030s; the CRS warns that real scientific and
engineering hurdles remain \citep{crs2026fusion}. Designing safety cases for
first-of-a-kind systems with limited operating experience is itself a challenge;
technology-agnostic Safety-in-Design methods are proposed
\citep{sidmethodology2026}.

% ----------------------------- RISK REGISTER -----------------------------
\section{Risk and uncertainty register}

The full register is in \texttt{deliverables/03-risk-uncertainty-register.md}.
The brief asks specifically for what is known, contested, and unknown.

\begin{description}[leftmargin=1.2em]
  \item[Known --- act confidently.] No chain reaction, meltdown, or long-lived
  high-level waste; tritium and activation are the real hazards; both the US and
  UK have committed to proportionate, fission-distinct regimes; a too-heavy
  regime carries real economic cost
  \citep{iterstafetybasis,ukfusionregfactsheet,nrc2023vote,crs2026fusion}.
  \item[Contested --- design for adaptability.] Whether the light regimes
  adequately bound tritium discharge and activated-waste handling; whether waste
  should be disposed of or recycled; how to build safety cases with little
  operating experience
  \citep{iaeatritiumharmonization,radwasterecycling2023,sidmethodology2026}.
  \item[Unknown --- resolve by monitoring.] Real activation inventories and waste
  volumes; the date and dominant technology of first grid connection; long-run
  performance of closed deuterium--tritium fuel cycles at scale
  \citep{tritiumbreeding2018,crs2026fusion,iaeatecdoc2049}.
\end{description}

The two highest-impact risks point in opposite directions: over-regulation (R1)
would stall a low-hazard industry, while under-regulation (R2) could miss a
hazard and trigger a confidence shock. Because several key parameters are
\emph{unknown} rather than merely contested, a fixed one-shot regime is poorly
suited to the problem.

% ----------------------------- OPTIONS -----------------------------
\section{Policy option appraisal}

Four options were scored against six weighted criteria using a Multi-Criteria
Decision Analysis (MCDA). The scoring is reproducible via
\texttt{verification/mcda\_score.py}; full pros, cons, and second-order effects
are in \texttt{deliverables/04-policy-option-appraisal.md}.

\begin{center}
\small
\begin{tabularx}{\linewidth}{@{}lXcccc@{}}
\toprule
\textbf{Criterion (weight)} & & \textbf{A} & \textbf{B} & \textbf{C} & \textbf{D} \\
\midrule
Safety (30)            & & 9 & 5 & 8 & 7 \\
Proportionality (20)   & & 2 & 9 & 8 & 6 \\
Innovation / speed (15)& & 2 & 9 & 7 & 4 \\
Adaptability (15)      & & 4 & 4 & 9 & 6 \\
Public trust (10)      & & 7 & 4 & 8 & 7 \\
Int'l coherence (10)   & & 5 & 4 & 7 & 9 \\
\midrule
\textbf{Weighted /100} & & \textbf{5.20} & \textbf{6.05} & \textbf{7.90} & \textbf{6.40} \\
\bottomrule
\end{tabularx}
\end{center}

\begin{description}[leftmargin=1.2em]
  \item[A --- Full fission-equivalent licensing.] Maximum assurance, but
  disproportionate, slow, and costly; risks driving developers offshore and
  entrenching the misleading ``fusion = fission'' framing.
  \item[B --- Light byproduct-only regime.] Fast and well-matched to the dominant
  hazard \citep{nrc2026proposedrule}, but thin on facility oversight and waste
  routes; vulnerable to a confidence shock \citep{iaeatritiumharmonization}.
  \item[C --- Proportionate, adaptive, review-gated framework (recommended).]
  Keeps the light byproduct core and adds environmental discharge control, worker
  safety, waste-clearance routes \citep{radwasterecycling2023}, and mandatory
  periodic review tied to new data and IAEA milestones. Highest on
  safety-adjusted proportionality and adaptability.
  \item[D --- International-harmonisation-first.] Best long-run coherence
  \citep{agilenations2023}, but IAEA standards are nascent and waiting forfeits
  near-term certainty \citep{iaeafusionsafety,crs2026fusion}.
\end{description}

A sensitivity check (swinging each weight by $\pm5$) leaves Option C top in every
scenario tested.

% ----------------------------- RECOMMENDATION -----------------------------
\section{Recommendation and implementation}

\keybox{\textbf{Recommendation: adopt Option C} --- a proportionate but adaptive,
review-gated framework. It preserves the fission-distinct stance the US and UK
have already chosen, while engineering in the adaptability the evidence demands.}

Practical steps for a jurisdiction adopting Option C:

\begin{enumerate}[leftmargin=1.5em]
  \item \textbf{Keep the proportionate core.} Regulate the radioactive materials
  (tritium, activation products), not the machine as a reactor
  \citep{nrc2026proposedrule,ukfusionregfactsheet}.
  \item \textbf{Add explicit controls.} Set tritium discharge limits and
  monitoring, worker-dose limits, emergency planning, and authorised
  waste-clearance and recycling routes \citep{iaeatritiumharmonization,radwasterecycling2023}.
  \item \textbf{Build in review gates.} Mandate scheduled reviews (for example
  every 3--5 years) triggered by new activation and tritium data and by IAEA
  harmonisation milestones \citep{iaea2023meeting,tritiumbreeding2018}.
  \item \textbf{Invest in regulator capability.} Fund fusion-specific expertise
  in the lead regulators so oversight can scale with the industry
  \citep{ukenvagency2022,sidmethodology2026}.
  \item \textbf{Lead, don't just follow, on harmonisation.} Engage the IAEA and
  Agile Nations process to shape common standards rather than inherit them
  \citep{agilenations2023}.
\end{enumerate}

This path lets policymakers give developers the certainty they need now, while
keeping the regime honest about what is still unknown.

\subsection{Indicative implementation timeline}
Option C can be sequenced without delaying developers' early-2030s targets:

\begin{description}[leftmargin=1.4em,style=nextline]
  \item[Year 0--1.] Adopt the proportionate byproduct / non-site-licence core in
  primary rules; publish interim tritium-discharge and worker-dose limits
  \citep{nrc2026proposedrule,ukfusionregfactsheet}.
  \item[Year 1--2.] Stand up fusion-specific capability in the lead regulators;
  authorise initial waste-clearance and recycling routes
  \citep{ukenvagency2022,radwasterecycling2023}.
  \item[Year 2--3.] Align reporting with IAEA tritium-release modelling and Agile
  Nations recommendations \citep{iaeatritiumharmonization,agilenations2023}.
  \item[Year 3--5 (first review gate).] Re-test discharge limits and waste
  classifications against the first commercial-scale activation data
  \citep{tritiumbreeding2018}; tighten or relax in proportion to evidence.
  \item[Every 3--5 years thereafter.] Repeat the review gate, keeping the regime
  matched to a maturing technology.
\end{description}

% ----------------------------- METHODOLOGY -----------------------------
\section{Methodology and notes on sources}

\textbf{Approach.} This briefing follows the UK Government Analysis Function
sequence: define the question, map stakeholders, synthesise evidence, separate
known/contested/unknown, appraise options, and recommend. Options were scored by
Multi-Criteria Decision Analysis with explicit, published weights, and the
scoring is reproducible from a script.

\textbf{Source base.} Twenty sources were used, spanning three classes: (i)
official and government material (NRC rulemaking and Federal Register notices, the
US Code of Federal Regulations, UK government strategy and legislation factsheets,
and a Congressional Research Service report); (ii) IAEA technical and convening
documents; and (iii) peer-reviewed literature (Philosophical Transactions of the
Royal Society A, the American Nuclear Society's \emph{Fusion Science and
Technology}, and the Springer \emph{Journal of Fusion Energy}). This mix exceeds
the requirement of fifteen peer-reviewed or official sources.

\textbf{Honesty and limitations.} This is a student-level briefing, not an
official regulatory instrument. Some commercialisation figures (for example,
total private investment) come from secondary reporting and are flagged as
indicative. Where activation inventories, waste volumes, or commercial dates are
genuinely unknown, the briefing says so rather than implying false precision.
Every substantive claim is tied to a numbered reference below.

% ----------------------------- REFERENCES -----------------------------
\newpage
\section*{References}
\addcontentsline{toc}{section}{References}
\renewcommand{\refname}{}
\vspace{-2.5em}
\bibliographystyle{unsrtnat}
\bibliography{../references}

\end{document}
